\documentclass[11pt]{article}
\usepackage[a4paper,margin=28mm]{geometry}
\usepackage{amsmath,amssymb,amsthm}
\usepackage[colorlinks=true,linkcolor=blue,urlcolor=blue,citecolor=blue]{hyperref}
\newtheorem*{theorem*}{Theorem}
\newtheorem*{corollary*}{Corollary}
\newenvironment{statement}[1]{\par\medskip\noindent\phantomsection\textbf{#1}\ }{\par\medskip}
\setlength{\emergencystretch}{3em}
\newcommand{\Q}{\mathbb Q}
\newcommand{\Z}{\mathbb Z}
\newcommand{\F}{\mathbb F}
\DeclareMathOperator{\End}{End}
\DeclareMathOperator{\Gal}{Gal}
\DeclareMathOperator{\GL}{GL}
\DeclareMathOperator{\PGL}{PGL}
\DeclareMathOperator{\Hom}{Hom}
\DeclareMathOperator{\Res}{Res}
\DeclareMathOperator{\Norm}{N}
\DeclareMathOperator{\Ad}{Ad}
\DeclareMathOperator{\tr}{tr}
\DeclareMathOperator{\ord}{ord}

\title{Rational \(3\)-torsion and the \(2\)-primary torsion of geometrically split abelian surfaces}
\author{}
\date{15 November 2026}

\begin{document}
\maketitle

\begin{theorem*}
Let \(A\) be an abelian surface over \(\Q\) that is simple over \(\Q\) and geometrically split. If \(A(\Q)\) has a
point of order \(3\), then \(A(\Q)\) has at most \(16\) points of \(2\)-power order:
\(\#A(\Q)[2^\infty]\le16\).
\end{theorem*}

\begin{corollary*}
No abelian surface over \(\Q\) that is simple over \(\Q\) and geometrically split has a rational torsion subgroup
of order divisible by \(96\). In particular
\href{https://genus-2-torsion.github.io/group.html?g=2.48}{\([2,48]\)},
\href{https://genus-2-torsion.github.io/group.html?g=2.2.24}{\([2,2,24]\)} and
\href{https://genus-2-torsion.github.io/group.html?g=2.2.2.12}{\([2,2,2,12]\)} do not occur as the rational
torsion subgroup of the Jacobian of a curve of genus \(2\) over \(\Q\) that is geometrically split and simple over
\(\Q\).
\end{corollary*}

Here ``geometrically split'' means that \(A_{\bar\Q}\) is isogenous to a product of two elliptic curves. We impose
no restriction on the degree of this isogeny, on the field of definition of the geometric endomorphisms or on the
integral endomorphism ring, and we assume no polarization. The bound \(16\) is attained, by the Jacobian of the
census curve 2.2.2.6.c, whose rational torsion subgroup is \([2,2,2,6]\) (\hyperref[rem-sharp]{Remark~2}). On
15 November 2026 the three groups of the Corollary are the only groups of the census whose order is divisible by
\(96\).

\emph{Outline.} We use the case distinction of \cite{P63}: \(A\) is of type (N), (CM) or (S). In type (CM) the
surface has potentially good reduction at \(5\), and no abelian surface over \(\F_5\) has a number of points
divisible by \(96\) (Section~\ref{five}). In the other types we move the torsion to one elliptic curve: in the
situations of Sections~\ref{quadratic} and~\ref{v4}, where \(F\) is a number field and \(A_F\) contains an
elliptic curve defined over \(F\), \(32\) rational points of \(2\)-power order and a rational point of order \(3\) give an elliptic curve
over \(F\) with a subgroup \(\Z/2\times\Z/12\) of \(F\)-rational points (Section~\ref{moving}). Over a quadratic field such curves exist, but each of them is the base
change of an elliptic curve over \(\Q\) (\hyperref[lemma-basechange]{Lemma~3}), and this contradicts the simplicity
of \(A\) over \(\Q\) (Section~\ref{quadratic}). This settles type (N), and type (S) when the field \(L\) of
definition of the geometric endomorphisms is quadratic. In type (S), the Frobenius polynomials at primes whose
Frobenius elements have order \(3\), \(4\) or \(6\) in \(\Gamma=\Gal(L/\Q)\) show that \(32\) rational points of
\(2\)-power order force \(\Gamma\cong C_2\) or \(\Gamma\cong V_4\) (Section~\ref{frobenius}). If
\(\Gamma\cong V_4\), the Frobenius polynomials at the involutions give \(\Q(i)\subset L\) from the \(2\)-primary
torsion and \(\Q(\sqrt{-3})\subset L\) from the point of order \(3\); so \(L=\Q(\zeta_{12})\), and no elliptic
curve over \(\Q(\zeta_{12})\) has a subgroup \(\Z/2\times\Z/12\) of rational points (Section~\ref{v4}).

\section{Conventions and inputs}\label{inputs}

All number fields lie in a fixed algebraic closure \(\bar\Q\), and \(G_k=\Gal(\bar\Q/k)\). For a prime \(\ell\),
\(\chi_\ell\) is the \(\ell\)-adic cyclotomic character, \(T_\ell\) the Tate module and
\(V_\ell=T_\ell\otimes\Q_\ell\). For an abelian variety \(X\), \(\End^0=\End\otimes\Q\) and
\(\Hom^0=\Hom\otimes\Q\). For a variety \(Y\) over \(\bar\Q\) and \(g\in G_{\Q}\), \(Y^g\) is the conjugate
variety, \(g\colon Y(\bar\Q)\to Y^g(\bar\Q)\) the induced bijection, and a morphism \(f\colon Y\to Z\) has the
conjugate \(g(f)\colon Y^g\to Z^g\), characterized by \(g(f)(g(P))=g(f(P))\) \cite[Section~1]{P2214}; so \(g\) acts
on the geometric endomorphisms of an abelian variety over \(\Q\) by \(g(\varphi)=g\circ\varphi\circ g^{-1}\). For
a quadratic field \(K\), \(\sigma\) is its nontrivial automorphism and \(\tilde\sigma\in G_{\Q}\) a lift of
\(\sigma\). For an abelian surface \(A\) over \(\Q\) and an elliptic curve \(E\subset A_K\) defined over \(K\),
\(E^\sigma=\tilde\sigma(E)\subset A_K\) is its conjugate; it is defined over \(K\) and does not depend on the lift
\cite[Section~1]{P40}. An elliptic curve in an abelian surface is an abelian subvariety of dimension \(1\). We write
\(\zeta_m\) for a primitive \(m\)-th root of unity, \(\Q^{\times2}\) for the set of nonzero rational squares, and
\([n_1,\dots,n_r]\) for the group \(\Z/n_1\times\dots\times\Z/n_r\), as on the census. We use the following
results.

\begin{itemize}
\item \textbf{(I1) Reduction} \cite{ST}; parts (a)--(c), and in (d) the absolute value of the roots, the number of
points and its bound, are as stated in \cite[(I2)]{P612} and \cite[(I1)]{P2210}. Let \(K'\) be a
finite extension of \(\Q_p\) with residue field \(k\), \(X\) an abelian variety over \(K'\), and \(\ell\neq p\) a
prime. (a) Good reduction is detected by the inertia action on \(T_\ell(X)\) \cite[\S1, Theorem~1]{ST}; it is
preserved by \(K'\)-isogenies \cite[\S1, Corollary~2]{ST}. (b) If \(X\) has good reduction \(\widetilde X\) and
\(m\) is prime to \(p\), then \(X(K')[m]\) injects into \(\widetilde X(k)\) \cite[Lemma~2]{ST}. (c) If \(X\) has
potentially good reduction, then it has good reduction over a finite extension of \(K'\) with residue field \(k\)
\cite[\S2]{ST}. (d) (Weil.) For an abelian surface \(B\) over \(\F_q\) and a prime \(\ell\nmid q\), the
characteristic polynomial \(P\in\Z[X]\) of the Frobenius endomorphism on \(V_\ell(B)\) does not depend on
\(\ell\); its complex roots have absolute value \(\sqrt q\); \(P(0)=q^2\), the degree of the Frobenius
endomorphism; and \(\#B(\F_q)=P(1)\). As \(P\) has real coefficients and \(P(0)>0\), its roots form two pairs
\(\alpha_i,\bar\alpha_i=q/\alpha_i\) (\(i=1,2\)): the roots that are not real come in conjugate pairs, and the
real roots \(\sqrt q\) and \(-\sqrt q\) each have even multiplicity, because the number of real roots is even and
the product of all roots is positive. So
\[P(X)=(X^2-\beta_1X+q)(X^2-\beta_2X+q)\]
with the real numbers \(\beta_i=\alpha_i+\bar\alpha_i\) of absolute value at most \(2\sqrt q\). Expanding,
\(P(X)=X^4+a_1X^3+a_2X^2+qa_1X+q^2\) with \(a_1=-(\beta_1+\beta_2)\) and \(a_2=2q+\beta_1\beta_2\), so the
polynomial \(h(y)=(y-\beta_1)(y-\beta_2)=y^2+a_1y+a_2-2q\) has integer coefficients; and
\[\#B(\F_q)=P(1)=h(1+q)\le(1+\sqrt q)^4.\]
If an abelian surface \(X\) over a number field \(k\) has good reduction \(\widetilde X\) at a prime
\(\mathfrak p\nmid\ell\), then \(P\) is the characteristic polynomial on \(V_\ell(X)\) of every arithmetic
Frobenius element in \(G_k\) at a place of \(\bar k\) above \(\mathfrak p\). For a prime \(p\neq\ell\) and an arithmetic
Frobenius element \(g\in G_{\Q}\) at \(p\), \(\chi_\ell(g)=p\).

\item \textbf{(I2) The case distinction} \cite[(I3)]{P612}, from \cite[Sections~2--5]{P63}. A geometrically split
abelian surface \(A/\Q\) is of exactly one of the types \textbf{(N)} (\(A_{\bar\Q}\) is isogenous to
\(E_1\times E_2\) with \(E_1\), \(E_2\) not isogenous), \textbf{(CM)} (\(A_{\bar\Q}\) is isogenous to \(E^2\) with
\(E\) a CM curve) and \textbf{(S)} (\(A_{\bar\Q}\) is isogenous to \(E^2\) with \(E\) without CM). In type (S),
\(D=\End^0(A_{\bar\Q})\cong M_2(\Q)\); \(L\) is the smallest field over which all geometric endomorphisms are
defined, a finite Galois extension of \(\Q\); \(\Gamma=\Gal(L/\Q)\) acts faithfully on \(D\), the element \(g\)
acting by \(\Ad(R_g)\colon\varphi\mapsto R_g\varphi R_g^{-1}\) with \(R_g\in D^\times\) unique up to
\(\Q^\times\); so \(\Gamma\) embeds in \(\PGL_2(\Q)\). The group \(\Gamma\) is one of \(C_1\), \(C_2\), \(C_3\),
\(C_4\), \(C_6\), \(V_4\), \(D_6\), \(D_8\), \(D_{12}\) (\(D_{2m}\) is dihedral of order \(2m\), and
\(V_4=C_2\times C_2\)), and \(\Gamma\neq C_1\) if \(A\) is simple over \(\Q\). In type (S), every elliptic curve
in \(A_{\bar\Q}\) is defined over \(L\), any two of them are \(L\)-isogenous, and for an elliptic curve
\(E\subset A_L\) the quotient \(A_L/E\) is \(L\)-isogenous to \(E\) \cite[Lemma~3]{P63}. In type (N), if \(A\) is
simple over \(\Q\), then \(A_{\bar\Q}\) contains exactly two elliptic curves; they are defined over a quadratic
field \(K\) and exchanged by \(\tilde\sigma\) \cite[Section~3]{P63}.

\item \textbf{(I3)} \cite[Lemma~2]{P63} Let \(X\) be an abelian surface over a number field \(k\), \(C\subset X\)
an elliptic curve defined over \(k\), and \(Q\in X(k)\) of prime order \(q\). Then either \(Q\in C(k)\), or the
image of \(Q\) in \((X/C)(k)\) has order \(q\).

\item \textbf{(I4) Frobenius elements of order \(3\), \(4\), \(6\)} \cite[Lemma~2]{P612}. Let \(A/\Q\) be of type
(S), \(\gamma\in\Gamma\) of order \(n\in\{3,4,6\}\), and \(p\) a prime of good reduction for \(A\), unramified in
\(L\), whose Frobenius class in \(\Gamma\) contains \(\gamma\). Then, for a primitive cube root of unity \(\zeta\)
and suitable integers \(a\), \(u\), \(b\),
\[\begin{array}{ll}
n=3\colon & \#\widetilde A(\F_p)=\Norm\bigl((1+a)+(a+p)\zeta\bigr),\\[2pt]
n=4\colon & \#\widetilde A(\F_p)=(1-u)^2+(p-u)^2,\\[2pt]
n=6\colon & \#\widetilde A(\F_p)=\Norm\bigl((1+2b+p)+(b+p)\zeta\bigr),
\end{array}\]
where \(\Norm(x+y\zeta)=x^2-xy+y^2\) is the norm of \(\Q(\zeta)/\Q\).

\item \textbf{(I5) Weil pairing} \cite[(I7)]{P2214}. For an isogeny \(u\colon E\to E'\) of elliptic curves and the
Weil pairings \(e\) on \(V_\ell\), \(e(ux,uy)=\deg(u)\,e(x,y)\); this extends to quasi-isogenies, with
\(\deg(u/m)=\deg(u)/m^2\). Moreover \(e(gx,gy)=\chi_\ell(g)\,e(x,y)\) for \(g\in G_{\Q}\), where
\(g\colon V_\ell(E)\to V_\ell(E^g)\).

\item \textbf{(I6) A Chebotarev criterion.} \emph{Let \(L/\Q\) be a finite Galois extension,
\(\gamma\in\Gal(L/\Q)\), \(m\ge1\) and \(a\) an integer prime to \(m\). Suppose that \(p\equiv a\bmod m\) for all
but finitely many primes \(p\) that are unramified in \(L\) and whose Frobenius class in \(\Gal(L/\Q)\) contains
\(\gamma\). Then \(\zeta_m\in L\) and \(\gamma(\zeta_m)=\zeta_m^a\).} Indeed, let \(M=L(\zeta_m)\) and let
\(\tau\in\Gal(M/\Q)\) be an extension of \(\gamma\). By Chebotarev's density theorem there are infinitely many
primes \(p\), unramified in \(M\), whose Frobenius class in \(\Gal(M/\Q)\) contains \(\tau\). For such \(p\) the
Frobenius class in \(\Gal(L/\Q)\) contains \(\tau|_L=\gamma\), and \(\tau(\zeta_m)=\zeta_m^p\) if \(p\nmid m\).
With finitely many exceptions these primes satisfy \(p\nmid m\) and \(p\equiv a\bmod m\), so
\(\tau(\zeta_m)=\zeta_m^a\). Hence all extensions of \(\gamma\) to \(M\) agree
on \(L\) and on \(\Q(\zeta_m)\), so \(\gamma\) has only one extension to \(M\). The number of extensions is
\([M:L]\), so \(M=L\).

\item \textbf{(I7) Elliptic curves.} Let \(k\) be a field of characteristic \(0\). (a) An elliptic curve over
\(k\) with invariants \(c_4\), \(c_6\) is isomorphic over \(k\) to \(y^2=x^3-27c_4x-54c_6\), and its discriminant
is \(\Delta=(c_4^3-c_6^2)/1728\) \cite[Section~III.1]{Si1}. The curve \(y^2=x^3+ax+b\) has invariants
\(c_4=-48a\), \(c_6=-864b\), and it is isomorphic over \(k\) to \(y^2=x^3+w^4ax+w^6b\) for every \(w\in k^\times\),
by \((x,y)\mapsto(w^2x,w^3y)\). (b) Division polynomials \cite[Exercise~III.3.7]{Si1}: for a Weierstrass equation
with coefficients \(a_1,\dots,a_6\) and \(n\ge1\) there is \(\psi_n\in\Z[a_1,\dots,a_6,x,y]\) such that, for every
elliptic curve \(E\) given by such an equation and every point \(P\neq O\) of \(E\), \(nP=O\) if and only if
\(\psi_n(P)=0\). (By part~(f) of the exercise, the divisor of the function \(\psi_n\) on \(E\) is the sum of the
points \(T\) with \(nT=O\), minus \(n^2\) times \(O\); so \(\psi_n\) vanishes exactly at the points \(P\neq O\)
with \(nP=O\).) On the curve, \(\psi_n\) is equal to a polynomial \(g_n\) in \(x\) for odd \(n\), and
to \(\psi_2g_n\) with \(g_n\) a polynomial in \(x\) for even \(n\), where \(\psi_2=2y+a_1x+a_3\) and, on the curve,
\(\psi_2^2=F(x)=4x^3+b_2x^2+2b_4x+b_6\). Put \(f_n=g_n\) for odd \(n\) and \(f_n=Fg_n\) for even \(n\). Then
\(nP=O\) if and only if \(f_n(x(P))=0\). The roots of \(F\) are the \(x\)-coordinates of the points of order
\(2\). (c) If \(X\) is an elliptic curve over a number field \(F\) with good reduction at a prime \(\mathfrak p\)
of residue characteristic \(p\), then the points of \(X(F)\) of finite order prime to \(p\) inject into the group
of points of the reduction over the residue field \cite[Proposition~VII.3.1]{Si1}. (d) For an elliptic curve \(X\)
over \(\Q\) and a biquadratic field \(L=\Q(\sqrt{d_1},\sqrt{d_2})\), the rank of \(X(L)\) is the sum of the ranks
over \(\Q\) of the quadratic twists of \(X\) by \(1\), \(d_1\), \(d_2\) and \(d_1d_2\): the
\(\Q\)-vector space \(X(L)\otimes\Q\) is the direct sum of the eigenspaces of the four characters of
\(\Gal(L/\Q)\), and the eigenspace of the character with kernel \(\Gal(L/\Q(\sqrt d))\) is isomorphic to
\(X^{(d)}(\Q)\otimes\Q\), since the twist \(X^{(d)}\) is isomorphic to \(X\) over \(\Q(\sqrt d)\) by an isomorphism
that the nontrivial automorphism of \(\Q(\sqrt d)\) maps to its negative. (e) A CM elliptic curve over a number
field has an integral \(j\)-invariant \cite[Theorem~II.6.1]{Si2}, hence potentially good reduction at every prime
\cite[Proposition~VII.5.5]{Si1}.
\end{itemize}

\section{A point count at \(5\)}\label{five}

\begin{statement}{Lemma 1 (point count at \(5\)).}\label{lemma-five}
(a) No abelian surface over \(\F_5\) has exactly \(96\) points over \(\F_5\).
(b) Let \(A\) be an abelian surface over \(\Q\) with potentially good reduction at \(5\). Then \(A(\Q)\) has no
subgroup of order \(96\).
\end{statement}

\emph{Proof.} (a) Let \(B\) be an abelian surface over \(\F_5\). By (I1)(d), \(\#B(\F_5)=h(6)\) with
\(h(y)=(y-\beta_1)(y-\beta_2)=y^2-sy+t\), where \(s,t\in\Z\) and \(\beta_1,\beta_2\) are real with
\(|\beta_i|\le2\sqrt5\). Suppose \(h(6)=36-6s+t=96\). Then \(t=60+6s\). As the roots of \(h\) are real,
\(s^2\ge4t=240+24s\), that is, \((s-12)^2\ge384\); as \(19^2=361<384\), this gives \(|s-12|\ge20\), so
\(s\le-8\) or \(s\ge32\). But \(|s|=|\beta_1+\beta_2|\le4\sqrt5<9\). Hence \(s=-8\), \(t=12\) and
\(h(y)=(y+2)(y+6)\), with the root \(-6<-2\sqrt5\), a contradiction. (Part~A of the certificate lists all values
of \(h(6)\) for polynomials \(h\) of this kind; the largest is \(100\), and \(96\) is not among them.)

(b) By (I1)(c) there is a finite extension \(K'\) of \(\Q_5\) with residue field \(\F_5\) over which \(A\) has good
reduction \(\widetilde A\). A subgroup of \(A(\Q)\subseteq A(K')\) of order \(96\) has order prime to \(5\), so by
(I1)(b) it injects into \(\widetilde A(\F_5)\), and \(96\) divides \(\#\widetilde A(\F_5)\). By (I1)(d),
\(\#\widetilde A(\F_5)\le(1+\sqrt5)^4=56+24\sqrt5<110<192\). So \(\#\widetilde A(\F_5)=96\), which contradicts
(a). \qed

\medskip
The prime \(5\) is special here: there are abelian surfaces over \(\F_7\) with \(96\) points, for instance the
products of an elliptic curve with \(8\) points and one with \(12\) points (Part~A). Lemma~1 uses neither
simplicity nor a polarization.

\section{Elliptic curves with a subgroup \(\Z/2\times\Z/12\)}\label{curves}

Let \(k\) be a field of characteristic \(0\). For \(t\in k\), \(t\neq1\), put
\[c(t)=-\frac{t(2t-1)(3t^2-3t+1)}{(t-1)^3},\qquad
b(t)=\frac{t(2t-1)(2t^2-2t+1)(3t^2-3t+1)}{(t-1)^4},\]
\[E_t\colon\ y^2+(1-c(t))xy-b(t)y=x^3-b(t)x^2.\]
This is Kubert's parametrization of the elliptic curves with a point of order \(12\) \cite[Table~3]{Ku}, in
factored form (Part~B of the certificate). The discriminant of this equation is
\[\Delta(E_t)=\frac{t^{12}(2t-1)^6(6t^2-6t+1)(3t^2-3t+1)^4(2t^2-2t+1)^3}{(t-1)^{24}}\]
(Part~B). So \(E_t\) is an elliptic curve if and only if \(D(t)\neq0\), where
\[D(t)=t(2t-1)(6t^2-6t+1)(3t^2-3t+1)(2t^2-2t+1).\]
Let \(X\) be the smooth projective curve over \(\Q\) with the affine equation
\[s^2=q(t),\qquad q(t)=4(6t^2-6t+1)(2t^2-2t+1)=48t^4-96t^3+80t^2-32t+4.\]
The polynomial \(q\) has four distinct roots, so \(X\) has genus \(1\), and it has two points at infinity, which
are conjugate over \(\Q(\sqrt3)\) because \(48=3\cdot4^2\). Let \(\Sigma\subset X(\bar\Q)\) be the set consisting of
the two points at infinity and of the affine points \((t,s)\) with \(t=1\) or \(D(t)=0\). We call the points of
\(\Sigma\) the \emph{cusps} of \(X\). The set \(\Sigma\) is stable under \(G_{\Q}\). It has \(16\) elements: the
points \((0,\pm2)\) and \((1,\pm2)\); the two points with \(t=1/2\), where \(q=-1\); the two points with \(s=0\) and
\(6t^2-6t+1=0\); the two points with \(s=0\) and \(2t^2-2t+1=0\); the four points with \(3t^2-3t+1=0\), where
\(q=-4/3\); and the two points at infinity. All of them are defined over \(\Q(i,\sqrt3)=\Q(\zeta_{12})\).

\begin{statement}{Lemma 2 (the curve \(X\)).}\label{lemma-curve}
(a) Let \(k\) be a field of characteristic \(0\) and \(E\) an elliptic curve over \(k\) such that \(E(k)\) has a
subgroup isomorphic to \(\Z/2\times\Z/12\). Then there is a point \((t_0,s_0)\in X(k)\) that is not a cusp such
that \(E\) is isomorphic over \(k\) to \(E_{t_0}\).
(b) The curve \(X\) is isomorphic over \(\Q\) to the elliptic curve with Cremona label \(24a4\), by an isomorphism
that maps \((0,2)\) to the origin. Its rank over \(\Q\) is \(0\), and
\(X(\Q)=\{(0,2),(0,-2),(1,2),(1,-2)\}\) is cyclic of order \(4\); all four rational points are cusps.
(c) \(X(\Q(\zeta_{12}))=\Sigma\). Hence no elliptic curve over \(\Q(\zeta_{12})\) has a subgroup of rational points
isomorphic to \(\Z/2\times\Z/12\).
\end{statement}

\emph{Proof.} (a) The group \(E(k)\) contains a point \(P\) of order \(12\) and all of \(E[2]\).
\begin{enumerate}
\item \emph{Tate normal form.} Choose a Weierstrass equation of \(E\) over \(k\) in which \(P=(0,0)\); then
\(a_6=0\). The tangent line at \(P\) is \(a_3y=a_4x\), and \(a_3\neq0\) because \(P\) is not of order \(2\).
Replacing \(y\) by \(y+(a_4/a_3)x\) gives an equation with \(a_4=a_6=0\), in which the tangent line \(y=0\) at
\(P\) meets the curve again in \((-a_2,0)=-2P\). As \(3P\neq O\), \(a_2\neq0\). The substitution
\((x,y)\mapsto(w^2x,w^3y)\) with \(w=a_3/a_2\) makes the coefficients \(a_2\) and \(a_3\) equal. So \(E\) is
isomorphic over \(k\) to a curve
\[E(b,c)\colon\ y^2+(1-c)xy-by=x^3-bx^2,\qquad b,c\in k,\]
by an isomorphism that maps \(P\) to \((0,0)\).

\item \emph{The order of \((0,0)\).} The discriminant of \(E(b,c)\) is a polynomial in \(b\) and \(c\) that is
divisible by \(b\), so \(b\neq0\). By (I7)(b), \(n\cdot(0,0)=O\) if and only if \(f_n(0)=0\), and \(f_n(0)\) is
a polynomial in \(b\) and \(c\). Part~B of the certificate computes the polynomials \(f_n\) for \(n\le12\) from
the recursion of \cite[Exercise~III.3.7]{Si1}, compares them with Magma's division polynomials, and factors
\(f_n(0)\): the irreducible factors of \(f_4(0)\), \(f_5(0)\), \(f_6(0)\) and \(f_{12}(0)\) are
\[\{b,\ c\},\qquad\{b,\ b-c\},\qquad\{b,\ b-c^2-c\},\qquad\{b,\ c,\ b-c^2-c,\ \Phi\},\]
respectively, where
\[\Phi=b^4-\tfrac13b^3c^2-3b^3c+\tfrac{10}3b^2c^2+\tfrac13bc^4-\tfrac53bc^3+\tfrac13c^6+\tfrac13c^4.\]
As \((0,0)\) has exact order \(12\) on \(E(b,c)\), we get \(c\neq0\), \(b\neq c\), \(b\neq c^2+c\) and
\(\Phi(b,c)=0\).

\item \emph{The parameter.} Put
\[T(b,c)=\frac{b^2-bc-c^3}{c(b-c^2-c)},\qquad\text{so that}\qquad T(b,c)-1=\frac{(b-c)^2}{c(b-c^2-c)}.\]
In \(\Q(b,c)\), the rational functions \(b(T(b,c))-b\) and \(c(T(b,c))-c\) have numerators divisible by \(\Phi\),
and denominators whose irreducible factors are among \(b\), \(c\), \(b-c\), \(b-c^2-c\) (Part~B). Let \(\Lambda\) be the
ring obtained from \(\Q[b,c]\) by inverting \(b\), \(c\), \(b-c\) and \(b-c^2-c\). It contains \(T(b,c)\) and the
inverse of \(T(b,c)-1\), hence also \(b(T(b,c))\) and \(c(T(b,c))\), because the denominators of \(b(t)\) and
\(c(t)\) are powers of \(t-1\); and \(b(T(b,c))-b\) and \(c(T(b,c))-c\) are multiples of \(\Phi\) in \(\Lambda\).
Evaluation at our pair \((b,c)\) is a ring homomorphism from \(\Lambda\) to \(k\), and \(\Phi\) vanishes at our pair.
Hence the element \(t_0=T(b,c)\in k\) is defined, \(t_0\neq1\), and \((b,c)=(b(t_0),c(t_0))\). So \(E\)
is isomorphic over \(k\) to \(E_{t_0}\), and \(D(t_0)\neq0\) because \(E_{t_0}\) is an elliptic curve.

\item \emph{The \(2\)-torsion.} For the equation of \(E_t\) over \(\Q(t)\), the polynomial \(F\) of (I7)(b)
factors as \(F(x)=(x-e_1(t))(4x^2+B(t)x+C(t))\), where \(e_1(t)=t(3t^2-3t+1)/(t-1)\) is the \(x\)-coordinate of
\(6\cdot(0,0)\), and the discriminant of the quadratic factor is
\[B(t)^2-16C(t)=q(t)\,\kappa(t)^2,\qquad\kappa(t)=\frac{(2t-1)^3(2t^2-2t+1)}{2(t-1)^6}\]
(Part~C). The coefficients of the equation of \(E_t\) and the rational functions \(e_1\), \(B\), \(C\) and
\(\kappa\) have denominators that are powers of \(t-1\) (Part~C), so these identities hold for \(t=t_0\). As
\(D(t_0)\neq0\), we have \(\kappa(t_0)\neq0\). All points of order \(2\) of \(E_{t_0}\) are defined over \(k\), so
the three roots of \(F\) lie in \(k\); hence the quadratic factor has its roots in \(k\), its discriminant is a
square in \(k\), and therefore \(q(t_0)=s_0^2\) with \(s_0\in k\). The point \((t_0,s_0)\in X(k)\) is not a cusp,
since \(t_0\neq1\) and \(D(t_0)\neq0\).
\end{enumerate}

(b) Part~D of the certificate: the elliptic curve \((X,(0,2))\) has Cremona label \(24a4\), rank \(0\) and torsion
subgroup \(\Z/4\) over \(\Q\), and the four points listed are rational. They are cusps, as \(t=0\) or \(t=1\).

(c) Let \(L=\Q(\zeta_{12})=\Q(i,\sqrt3)\). By (I7)(d) the rank of \(X(L)\) is the sum of the ranks over \(\Q\) of
the twists of \(X\) by \(1\), \(-1\), \(-3\) and \(3\); these are the curves \(24a4\), \(48a4\), \(72a1\) and
\(144b1\), all of rank \(0\) (Parts~D and~F). So \(X(L)\) is finite. The primes \(13\) and \(37\) split completely
in \(L\), and the Weierstrass model of \(X\) used in the certificate has good reduction at both (Part~F). By (I7)(c), a point of order \(13\) in \(X(L)\) would give a point
of order \(13\) in \(X(\F_{37})\); but \(13\nmid\#X(\F_{37})\) (Part~F). So \(X(L)\) has order prime to \(13\), and
by (I7)(c) it injects into \(X(\F_{13})\), a group of order \(16\) (Part~F). As \(\Sigma\subseteq X(L)\) has
\(16\) elements, \(X(L)=\Sigma\). (Part~F also computes the \(16\) points of \(X(L)\) and checks that each of them
is a cusp.) The last assertion follows from (a). \qed

\begin{statement}{Lemma 3 (base change).}\label{lemma-basechange}
Let \(K\) be a quadratic field and \(E\) an elliptic curve over \(K\) such that \(E(K)\) has a subgroup isomorphic
to \(\Z/2\times\Z/12\). Then \(E\) is isomorphic over \(K\) to the base change \((E_0)_K\) of an elliptic curve
\(E_0\) over \(\Q\).
\end{statement}

\emph{Proof.} By \hyperref[lemma-curve]{Lemma~2}(a), \(E\) is isomorphic over \(K\) to \(E_{t_0}\) for a point
\(Q=(t_0,s_0)\in X(K)\) that is not a cusp. We use the group law of \(X\) with origin \(O=(0,2)\). The point
\(R=Q+\sigma(Q)\) is fixed by \(\sigma\), so \(R\in X(\Q)\), and \(\sigma(Q)=\iota_R(Q)\), where
\(\iota_R\colon X\to X\), \(P\mapsto R-P\), is an involution of \(X\) defined over \(\Q\). Let \(c_4\), \(c_6\)
and \(\Delta\) be the invariants and the discriminant of \(E_t\), which are elements of \(\Q(t)\subset\Q(X)\) whose
denominators are powers of \(t-1\) (Part~B).

\emph{The functions \(\lambda_R\).} For each of the four points \(R\in X(\Q)\), Part~E of the certificate computes
\(\iota_R\) and a function \(\lambda_R\in\Q(X)\) with
\[c_4\circ\iota_R=\lambda_R^4\,c_4,\qquad c_6\circ\iota_R=\lambda_R^6\,c_6,\qquad
\lambda_R\cdot(\lambda_R\circ\iota_R)=1.\]
Explicitly, with \(w_0=s/(4(6t^2-6t+1))\), \(l_0=(3t^2-3t+1)/(2t^3)\) and
\(l_1=(5t^2-5t+1)s/(4t^3(6t^2-6t+1))\):
\[\begin{array}{lll}
R=(0,2)\colon & t\circ\iota_R=\tfrac12-w_0, & \lambda_R=l_0-l_1;\\[2pt]
R=(0,-2)\colon & \iota_R(t,s)=(t,-s), & \lambda_R=1;\\[2pt]
R=(1,2)\colon & \iota_R(t,s)=(1-t,s), & \lambda_R=(t-1)^3/t^3;\\[2pt]
R=(1,-2)\colon & t\circ\iota_R=\tfrac12+w_0, & \lambda_R=l_0+l_1.
\end{array}\]
Consequently \(\Delta\circ\iota_R=\lambda_R^{12}\,\Delta\).

\emph{The value of \(\lambda_R\) at \(Q\).} At a point of \(X\) that is not a cusp, the functions \(c_4\), \(c_6\)
and \(\Delta\) are regular and \(\Delta\) does not vanish, by the formula for \(\Delta(E_t)\). As \(\Sigma\) is
stable under \(G_{\Q}\), the point \(\iota_R(Q)=\sigma(Q)\) is not a cusp either. So
\(\ord_Q(\Delta\circ\iota_R)=\ord_{\iota_R(Q)}(\Delta)=0\) and \(\ord_Q(\Delta)=0\), hence
\(\ord_Q(\lambda_R)=0\): the function \(\lambda_R\) is regular at \(Q\) and \(u=\lambda_R(Q)\in K^\times\).
The same holds at \(\sigma(Q)\). As \(c_4\), \(c_6\), \(\lambda_R\) and \(\iota_R\) are defined over \(\Q\),
\[\sigma\bigl(c_4(Q)\bigr)=c_4(\sigma(Q))=c_4(\iota_R(Q))=u^4c_4(Q),\qquad
\sigma\bigl(c_6(Q)\bigr)=u^6c_6(Q),\]
\[u\,\sigma(u)=\lambda_R(Q)\,\lambda_R(\sigma(Q))=\lambda_R(Q)\,\lambda_R(\iota_R(Q))=1.\]

\emph{Descent.} By Hilbert's Theorem~90 there is \(w\in K^\times\) with \(u=w/\sigma(w)\); explicitly \(w=1+u\) if
\(u\neq-1\), and \(w=\sqrt d\) if \(u=-1\) and \(K=\Q(\sqrt d)\). Then \(A_4=w^4c_4(Q)\) and \(A_6=w^6c_6(Q)\) are
fixed by \(\sigma\), so they are rational. Let \(E_0\) be the curve \(y^2=x^3-27A_4x-54A_6\) over \(\Q\). By
(I7)(a) its invariants are \((6w)^4c_4(Q)\) and \((6w)^6c_6(Q)\), and its discriminant is
\(6^{12}w^{12}\Delta(Q)\neq0\); so \(E_0\) is an elliptic curve, and \((E_0)_K\) is isomorphic over \(K\) to
\(y^2=x^3-27c_4(Q)x-54c_6(Q)\), that is, to \(E_{t_0}\), and so to \(E\). \qed

\medskip
Part~G of the certificate carries out this descent for one quadratic point \(Q\) for each of the four points
\(R\). It also checks the statement of Lemma~3 in a way that does not use the involutions \(\iota_R\), for twelve
points of \(X\) over \(\Q(\sqrt{13})\), a field over which \(X\) has positive rank
(\hyperref[rem-bn]{Remark~5}).

\section{Moving the torsion to one elliptic curve}\label{moving}

\begin{statement}{Lemma 4 (splitting the torsion).}\label{lemma-split}
Let \(F\) be a number field, \(X\) an abelian surface over \(F\), \(C\subset X\) an elliptic curve defined over
\(F\), and \(\pi\colon X\to Y=X/C\) the quotient map, so that \(Y\) is an elliptic curve over \(F\).
(a) Let \(G\subseteq X(F)\) be a finite \(2\)-group of order at least \(32\). Then \(G\cap C(F)\) or
\(\pi(G)\subseteq Y(F)\) has a subgroup isomorphic to \(\Z/8\) or to \(\Z/2\times\Z/4\).
(b) Let \(P\in X(F)\) be a point of order \(3\). Then \(P\in C(F)\), or \(\pi(P)\in Y(F)\) has order \(3\).
\end{statement}

\emph{Proof.} (a) The kernel of \(\pi\) on \(G\) is \(G\cap C(\bar\Q)=G\cap C(F)\), so
\(\#G=\#(G\cap C(F))\cdot\#\pi(G)\), and one of the two groups has order at least \(8\). It is a finite
\(2\)-group of rational points of an elliptic curve, so it is isomorphic to \(\Z/2^\alpha\times\Z/2^\beta\) with
\(\alpha\le\beta\) and \(\alpha+\beta\ge3\). If \(\alpha=0\), it contains \(\Z/8\); if \(\alpha\ge1\), then
\(\beta\ge2\) and it contains \(\Z/2\times\Z/4\). (b) This is (I3). \qed

\begin{statement}{Lemma 5 (combining).}\label{lemma-combine}
Let \(F\) be a number field and \(E_a\), \(E_b\) elliptic curves over \(F\) that are \(F\)-isogenous or equal.
Suppose that \(E_a(F)\) has a subgroup \(T\) isomorphic to \(\Z/8\) or to \(\Z/2\times\Z/4\), and that \(E_b(F)\)
has a point \(P\) of order \(3\). Then there is an elliptic curve \(E_3\) over \(F\), \(F\)-isogenous to \(E_a\),
such that \(E_3(F)\) has a subgroup isomorphic to \(\Z/2\times\Z/12\).
\end{statement}

\emph{Proof.} Let \(f\colon E_a\to E_b\) be an \(F\)-isogeny (the identity if \(E_a=E_b\)). Let \(H\) be the
\(3\)-primary part of \(\ker f\); it is a characteristic subgroup of the \(G_F\)-stable group \(\ker f\), hence
\(G_F\)-stable. Put \(E'=E_a/H\), let \(a\colon E_a\to E'\) be the quotient map and write \(f=b\circ a\) with
\(b\colon E'\to E_b\) defined over \(F\); the degree of \(a\) is a power of \(3\) and the degree of \(b\) is prime
to \(3\). The map \(a\) is injective on the points of \(2\)-power order, so \(a(T)\subseteq E'(F)\) is isomorphic
to \(T\). Since \(b\circ\hat b=[\deg b]\) is injective on \(E_b[3]\), so is the dual isogeny
\(\hat b\colon E_b\to E'\), which is defined over \(F\); so \(\hat b(P)\in E'(F)\) has order \(3\). Hence \(E'(F)\)
contains \(a(T)\times\langle\hat b(P)\rangle\), which is isomorphic to \(\Z/24\) or to \(\Z/2\times\Z/12\). In the
second case put \(E_3=E'\).

In the first case let \(Q\in E'(F)\) be a point of order \(24\), put \(E_3=E'/\langle12Q\rangle\), and let
\(\varphi\colon E'\to E_3\) be the quotient map; it is defined over \(F\). Then \(\varphi(Q)\in E_3(F)\) has order
\(12\), because \(n\varphi(Q)=O\) if and only if \(nQ\in\{O,12Q\}\). Write \(E'[2]=\{O,12Q,T_1,T_2\}\). The group
\(G_F\) fixes \(O\) and \(12Q\), so it preserves \(\{T_1,T_2\}\); and \(T_2=T_1+12Q\), so
\(\varphi(T_1)=\varphi(T_2)\). Hence \(W=\varphi(T_1)\) is fixed by \(G_F\): \(W\in E_3(F)\). It has order \(2\),
as \(T_1\notin\ker\varphi\). Moreover \(W\neq6\varphi(Q)\): otherwise \(T_1-6Q\in\ker\varphi=\{O,12Q\}\), but
\(2(T_1-6Q)=12Q\neq O\), so \(T_1-6Q\) has order \(4\). As \(6\varphi(Q)\) is the only point of order \(2\) in
\(\langle\varphi(Q)\rangle\), the group \(E_3(F)\) contains
\(\langle W\rangle\times\langle\varphi(Q)\rangle\cong\Z/2\times\Z/12\). In both cases \(E_3\) is \(F\)-isogenous
to \(E'\), hence to \(E_a\). \qed

\section{The quadratic case}\label{quadratic}

\begin{statement}{Theorem 6 (quadratic case).}\label{thm-quadratic}
Let \(A\) be an abelian surface over \(\Q\) that is simple over \(\Q\), and \(K\) a quadratic field such that
\(A_K\) contains an elliptic curve \(E\) defined over \(K\). If \(A(\Q)\) has a point of order \(3\), then
\(\#A(\Q)[2^\infty]\le16\).
\end{statement}

\emph{Proof.} Suppose that \(A(\Q)\) has a point \(P_3\) of order \(3\) and that \(G=A(\Q)[2^\infty]\) has order
at least \(32\).

\emph{Step 1: the curves \(E\) and \(Y\).} An elliptic curve \(E\subset A_K\) with \(E^\sigma=E\) is stable under
\(G_{\Q}\), hence an elliptic curve in \(A\) defined over \(\Q\) \cite[Section~2]{P2214}; as \(A\) is simple over
\(\Q\), \(E^\sigma\neq E\). So \(E\) and \(E^\sigma\) are distinct elliptic curves in \(A_K\), and \(E\cap E^\sigma\)
is finite. Let \(\pi\colon A_K\to Y=A_K/E^\sigma\) be the quotient map, defined over \(K\). Its restriction to
\(E\) has the finite kernel \(E\cap E^\sigma\), so it is a \(K\)-isogeny \(E\to Y\). If \(P\in A(\Q)\) lies in
\(E^\sigma\), then \(P=\tilde\sigma(P)\in\tilde\sigma(E^\sigma)=\tilde\sigma^2(E)=E\), as \(\tilde\sigma^2\in G_K\);
so \(P\in E(\bar\Q)\cap A(K)=E(K)\). Thus \(A(\Q)\cap E^\sigma(K)\subseteq E(K)\).

\emph{Step 2: an elliptic curve over \(K\) with a subgroup \(\Z/2\times\Z/12\).} We apply
\hyperref[lemma-split]{Lemma~4} to \(X=A_K\), \(C=E^\sigma\) and \(G\). By (a) and Step~1, one of \(E(K)\),
\(Y(K)\) has a subgroup \(T\) isomorphic to \(\Z/8\) or to \(\Z/2\times\Z/4\); by (b) and Step~1, one of \(E(K)\),
\(Y(K)\) has a point of order \(3\). As \(E\) and \(Y\) are \(K\)-isogenous,
\hyperref[lemma-combine]{Lemma~5} gives an elliptic curve \(E_3\) over \(K\), \(K\)-isogenous to \(E\), such that
\(E_3(K)\) has a subgroup isomorphic to \(\Z/2\times\Z/12\).

\emph{Step 3: descent.} By \hyperref[lemma-basechange]{Lemma~3}, \(E_3\) is isomorphic over \(K\) to \((E_0)_K\)
for an elliptic curve \(E_0\) over \(\Q\). So there is a \(K\)-isogeny \(\psi\colon(E_0)_K\to E\). Let
\(h_1\colon(E_0)_K\to A_K\) be the composition of \(\psi\) with the inclusion of \(E\); it is a homomorphism
defined over \(K\) with image \(E\). As \(E_0\) and \(A\) are defined over \(\Q\), the conjugate
\(\tilde\sigma(h_1)\colon(E_0)_K\to A_K\) is a homomorphism defined over \(K\), with image
\(\tilde\sigma(E)=E^\sigma\), and \(\tilde\sigma(\tilde\sigma(h_1))=h_1\) because \(\tilde\sigma^2\in G_K\). Hence
\(\Psi=h_1+\tilde\sigma(h_1)\) is fixed by \(G_K\) and by \(\tilde\sigma\), so by \(G_{\Q}\): it is a homomorphism
\(E_0\to A\) defined over \(\Q\). If \(\Psi=0\), then \(h_1=-\tilde\sigma(h_1)\), and the images \(E\) and
\(E^\sigma\) of these two maps would be equal. So \(\Psi\neq0\), and the image of \(\Psi\) is an elliptic curve in
\(A\) defined over \(\Q\). This contradicts the simplicity of \(A\) over \(\Q\). \qed

\section{Type (S): Frobenius elements}\label{frobenius}

In this section and the next, \(A/\Q\) is of type (S) and simple over \(\Q\), with \(D\), \(L\) and \(\Gamma\) as
in (I2).

\begin{statement}{Lemma 7 (Frobenius at involutions).}\label{lemma-inv}
Let \(\gamma\in\Gamma\) be an element of order \(2\).
(a) Every \(R\in D^\times\) with \(\gamma=\Ad(R)\) satisfies \(\tr R=0\) and \(R^2\in\Q^\times\), and the class of
\(R^2\) in \(\Q^\times/\Q^{\times2}\) depends only on \(\gamma\). We choose \(R\) such that \(c=R^2\) is a
squarefree integer, and let \(\varepsilon\in\{1,-1\}\) be the sign of \(c\).
(b) Let \(p\) be a prime of good reduction for \(A\), unramified in \(L\), whose Frobenius class in \(\Gamma\)
contains \(\gamma\). Then the characteristic polynomial \(P\) of Frobenius of \(\widetilde A/\F_p\) and the number
of points are
\[P(X)=(X^2+\varepsilon p)^2-cw^2X^2,\qquad\#\widetilde A(\F_p)=(1+\varepsilon p)^2-cw^2,\]
for an integer \(w\ge0\).
\end{statement}

\emph{Proof.} (a) As \(\gamma^2=1\), \(\Ad(R^2)\) is the identity, so \(R^2\) lies in the centre \(\Q\) of
\(D\cong M_2(\Q)\). As \(\Gamma\) acts faithfully, \(R\) is not a scalar. So \(X^2-R^2\) is the minimal polynomial
of \(R\), hence its characteristic polynomial, and \(\tr R=0\). The elements inducing \(\gamma\) are the multiples
\(rR\), \(r\in\Q^\times\), and \((rR)^2=r^2R^2\).

(b) Let \(\ell\neq p\) be a prime, \(\mathfrak P\) a prime of \(L\) above \(p\) with Frobenius element \(\gamma\),
and \(g\in G_{\Q}\) an arithmetic Frobenius element at a place of \(\bar\Q\) above \(\mathfrak P\). Then
\(g|_L=\gamma\), \(\chi_\ell(g)=p\), and \(P\) is the characteristic polynomial of \(g\) on \(V_\ell(A)\), by
(I1)(d).
\begin{enumerate}
\item \emph{The multiplicity space} \cite[Lemma~3, Step~1]{P2214}. Choose an elliptic curve \(E\subset A_L\); it
is defined over \(L\) by (I2). Put \(H=\Hom^0(E_{\bar\Q},A_{\bar\Q})\), a \(\Q\)-vector space of dimension \(2\).
The algebra \(D\) acts on \(H\) by composition, faithfully because \(D\) is simple and the action is unital; so
\(D=\End_{\Q}(H)\), as \(\dim_{\Q}H=2\). The evaluation map
\[\mathrm{ev}\colon V_\ell(E)\otimes_{\Q}H\to V_\ell(A),\qquad v\otimes h\mapsto h(v),\]
is an isomorphism: for a basis \(h_1,h_2\) of \(H\), a nonzero integer multiple of
\((h_1,h_2)\colon E^2\to A_{\bar\Q}\) is a homomorphism, and it is an isogeny (otherwise its image would be an
elliptic curve \(E_1\), and \(h_1,h_2\) would be linearly dependent in the one-dimensional space
\(\Hom^0(E,E_1)\)); \(\mathrm{ev}\) is the map it induces, up to that integer.

\item \emph{The action of \(g\)} \cite[Lemma~3, Step~2]{P2214}. The curve \(E^g\subset A_{\bar\Q}\) is defined over
\(L\) and \(L\)-isogenous to \(E\) by (I2); choose an invertible \(u\in\Hom^0(E,E^g)\). Define
\(B\colon V_\ell(E)\to V_\ell(E)\) by \(B(v)=u^{-1}(g(v))\) and \(C\colon H\to H\) by \(C(h)=g(h)\circ u\). For
\(v\in V_\ell(E)\) and \(h\in H\),
\[g\bigl(h(v)\bigr)=g(h)\bigl(g(v)\bigr)=g(h)\bigl(u(Bv)\bigr)=C(h)(Bv),\]
so \(g\circ\mathrm{ev}=\mathrm{ev}\circ(B\otimes C)\). The map \(C\) is \(\Q\)-linear and invertible, with inverse
\(h\mapsto g^{-1}(h\circ u^{-1})\). For \(\varphi\in D\) we have \(C(\varphi\circ h)=g(\varphi)\circ C(h)\); so
\(C\in\GL(H)=D^\times\) satisfies \(g(\varphi)=C\varphi C^{-1}\). By (I5),
\(e(Bv,Bv')=\deg(u)^{-1}e(g(v),g(v'))=\deg(u)^{-1}\chi_\ell(g)\,e(v,v')\) for \(v,v'\in V_\ell(E)\) and the Weil
pairing \(e\) of \(E\), so
\[\delta:=\det B=p/\deg(u)\]
is a positive rational number.

\item \emph{The shape of \(P\).} The action of \(g\) on \(D\) is that of \(g|_L=\gamma\). So \(\Ad(C)=\Ad(R)\),
and \(C=rR\) with \(r\in\Q^\times\). Hence \(C^2=r^2c\) and \(\tr C=0\): the eigenvalues of \(C\) are
\(\pm\rho\) with \(\rho^2=r^2c\). Let \(\beta_1,\beta_2\) be the eigenvalues of \(B\), so
\(\beta_1\beta_2=\delta\). By Step~2 the eigenvalues of \(g\) on \(V_\ell(A)\) are \(\pm\rho\beta_1\),
\(\pm\rho\beta_2\), so
\[P(X)=(X^2-r^2c\beta_1^2)(X^2-r^2c\beta_2^2)=(X^2+r^2c\delta)^2-c\,\theta^2X^2\]
with \(\theta=r(\beta_1+\beta_2)=r\tr B\in\Q_\ell\) (the identity is checked in Part~H). As \(P(0)=p^2\), \(r^2c\delta=\pm p\); as \(r^2\delta>0\), the sign is that of
\(c\). So \(P(X)=(X^2+\varepsilon p)^2-c\,\theta^2X^2\).

\item \emph{\(\theta^2\) is the square of an integer.} Write \(P(X)=X^4+a_1X^3+a_2X^2+\dots\) with \(a_i\in\Z\); these
coefficients do not depend on \(\ell\). By Step~3, \(a_1=0\) and \(m:=(2\varepsilon p-a_2)/c=\theta^2\). So the rational number
\(m\) is a square in \(\Q_\ell\) for every prime \(\ell\neq p\). Then \(m\) is a square in \(\Q\): otherwise
\(\Q(\sqrt m)\) would be a quadratic field, and \(m\) would not be a square in \(\Q_\ell\) for the infinitely many
primes \(\ell\) that are inert in it. Write \(m=w^2\) with \(w\in\Q\), \(w\ge0\). As \(cw^2=2\varepsilon p-a_2\) is an integer
and \(c\) is squarefree, \(w\) is an integer. Finally \(\#\widetilde A(\F_p)=P(1)=(1+\varepsilon p)^2-cw^2\). \qed
\end{enumerate}

\begin{statement}{Proposition 8 (rotations).}\label{prop-rot}
If \(\#A(\Q)[2^\infty]\ge32\), then \(\Gamma\) has no element of order \(3\), \(4\) or \(6\); so
\(\Gamma\cong C_2\) or \(\Gamma\cong V_4\).
\end{statement}

\emph{Proof.} Suppose that \(\gamma\in\Gamma\) has order \(n\in\{3,4,6\}\). Let \(p\) be an odd prime of good
reduction for \(A\), unramified in \(L\), whose Frobenius class in \(\Gamma\) contains \(\gamma\). By (I1)(b),
\(A(\Q)[2^\infty]\) injects into \(\widetilde A(\F_p)\), so \(32\) divides \(\#\widetilde A(\F_p)\), which is given
by (I4). We claim that \(p\equiv1\bmod8\).

For \(n=4\): \(x^2+y^2\equiv0\bmod32\) with \(x=1-u\), \(y=p-u\). A sum of two integer squares that is divisible
by \(4\) has both terms even; applying this twice, \(x=4x''\) and \(y=4y''\) with \(x''^2+y''^2\) even. So
\(x''\equiv y''\bmod2\), and \(8\) divides \(y-x=p-1\).
For \(n=3\) and \(n=6\): the prime \(2\) is inert in \(\Z[\zeta]\), so the \(2\)-adic valuation of
\(\Norm(x+y\zeta)\) is twice the largest \(v\) with \(2^v\mid x\) and \(2^v\mid y\). If \(32\) divides the norm,
then \(v\ge3\): \(8\mid x\) and \(8\mid y\). For \(n=3\), \(x=1+a\) and \(y=a+p\), so \(8\) divides \(y-x=p-1\).
For \(n=6\), \(x=1+2b+p\) and \(y=b+p\), so \(8\) divides \(x-y=1+b\), and \(p\equiv-b\equiv1\bmod8\).
(Part~H checks the three cases over all residues modulo \(32\).)

So \(p\equiv1\bmod8\) for all but finitely many primes \(p\) that are unramified in \(L\) and whose Frobenius class
contains \(\gamma\). By (I6) with \(m=8\) and \(a=1\), \(\zeta_8\in L\) and \(\gamma\) fixes \(\zeta_8\). Hence
all elements of \(\Gamma\) of order \(3\), \(4\) or \(6\) lie in \(\Gal(L/\Q(\zeta_8))\), a subgroup of index
\(4\). But in each of the groups \(C_3\), \(C_4\), \(C_6\), \(D_6\), \(D_8\), \(D_{12}\) the elements of order
\(3\), \(4\) or \(6\) generate a subgroup of index at most \(2\) (the group itself, or its cyclic subgroup of
index \(2\)). This is a contradiction. By (I2), \(\Gamma\) is then \(C_1\), \(C_2\) or \(V_4\), and
\(\Gamma\neq C_1\) because \(A\) is simple over \(\Q\). \qed

\section{Type (S) with \(\Gamma\cong V_4\)}\label{v4}

\begin{statement}{Proposition 9 (\(\Gamma\cong V_4\)).}\label{prop-v4}
Suppose that \(\Gamma\cong V_4\).
(a) If \(\#A(\Q)[2^\infty]\ge32\), then \(\Q(i)\subset L\).
(b) If \(A(\Q)\) has a point of order \(3\), then \(\Q(\sqrt{-3})\subset L\).
(c) If \(A(\Q)\) has a point of order \(3\), then \(\#A(\Q)[2^\infty]\le16\).
\end{statement}

\emph{Proof.} Let \(\gamma_1,\gamma_2,\gamma_3\) be the elements of order \(2\) of \(\Gamma\), and let \(R_i\),
\(c_i=R_i^2\) and \(\varepsilon_i\) be as in \hyperref[lemma-inv]{Lemma~7}(a) for \(\gamma_i\).

\emph{The classes \(c_i\).} As \(\gamma_1\gamma_2=\gamma_2\gamma_1\), we have \(R_1R_2=\mu R_2R_1\) with
\(\mu\in\Q^\times\), and \(\mu^2=1\) by taking determinants. If \(\mu=1\), then \(R_2\) lies in the centralizer
\(\Q+\Q R_1\) of the non-scalar matrix \(R_1\); as \(\tr R_2=\tr R_1=0\), \(R_2\in\Q R_1\) and
\(\gamma_2=\gamma_1\), which is false. So \(R_1R_2=-R_2R_1\) and \((R_1R_2)^2=-R_1^2R_2^2=-c_1c_2\). As \(R_1R_2\)
induces \(\gamma_1\gamma_2=\gamma_3\), Lemma~7(a) gives \(c_3\in-c_1c_2\,\Q^{\times2}\). Hence
\[c_1c_2c_3\in-\Q^{\times2}.\]

(a) As \(-1\) is not a square in \(\Q_2\), some \(c_i\) is not a square in \(\Q_2\). This \(c_i\) is a
squarefree integer, so \(c_i\not\equiv0\bmod4\), and \(c_i\not\equiv1\bmod8\) because the integers congruent to
\(1\) modulo \(8\) are squares in \(\Q_2\). We use: \emph{if \(c\not\equiv0\bmod4\), \(c\not\equiv1\bmod8\) and
\(8\mid x^2-cy^2\), then \(x\) and \(y\) are even.} Indeed, if \(y\) were odd, then \(x^2\equiv c\bmod8\), but the
squares modulo \(8\) are \(0\), \(1\) and \(4\); so \(y\) is even, \(4\mid x^2\), and \(x\) is even. Now let \(p\)
be an odd prime of good reduction for \(A\), unramified in \(L\), whose Frobenius class contains \(\gamma_i\). By
(I1)(b) and Lemma~7(b), \(32\) divides \(x^2-c_iw^2\) with \(x=1+\varepsilon_ip\). Applying the statement in italics to
\((x,w)\) and then to \((x/2,w/2)\), we get \(4\mid x\), that is, \(p\equiv-\varepsilon_i\bmod4\) (Part~H checks this over
all residues modulo \(32\)). By (I6) with \(m=4\) and \(a=-\varepsilon_i\), \(i\in L\).

(b) As \(-1\) is not a square in \(\Q_3\), some \(c_j\) is not a square in \(\Q_3\), so \(c_j\not\equiv1\bmod3\).
Let \(p\neq3\) be a prime of good reduction for \(A\), unramified in \(L\), whose Frobenius class contains
\(\gamma_j\). By (I1)(b) and Lemma~7(b), \(3\) divides \(x^2-c_jw^2\) with \(x=1+\varepsilon_jp\). If \(c_j\equiv0\bmod3\),
then \(3\mid x\). If \(c_j\equiv2\bmod3\), then \(x^2\equiv2w^2\bmod3\), and \(3\mid x\) because \(2\) is not a
square modulo \(3\). So \(p\equiv-\varepsilon_j\bmod3\), and by (I6) with \(m=3\) and \(a=-\varepsilon_j\), \(\zeta_3\in L\).

(c) Suppose that \(A(\Q)\) has a point of order \(3\) and that \(G=A(\Q)[2^\infty]\) has order at least \(32\).
By (a) and (b), \(L\) contains \(\Q(i,\sqrt{-3})=\Q(\zeta_{12})\); as \([L:\Q]=4\), \(L=\Q(\zeta_{12})\). Choose an
elliptic curve \(E\subset A_L\); it is defined over \(L\), and \(Y=A_L/E\) is \(L\)-isogenous to \(E\), by (I2).
By \hyperref[lemma-split]{Lemma~4} for \(X=A_L\) and \(C=E\), one of \(E(L)\), \(Y(L)\) has a subgroup
isomorphic to \(\Z/8\) or to \(\Z/2\times\Z/4\), and one of them has a point of order \(3\). By
\hyperref[lemma-combine]{Lemma~5} there is an elliptic curve over \(L=\Q(\zeta_{12})\) with a subgroup of
rational points isomorphic to \(\Z/2\times\Z/12\). This contradicts
\hyperref[lemma-curve]{Lemma~2}(c). \qed

\section{Proof of the Theorem and consequences for the census}\label{census}

\emph{Proof of the Theorem.} Let \(A/\Q\) be simple over \(\Q\) and geometrically split, with a rational point of
order \(3\), and suppose that \(\#A(\Q)[2^\infty]\ge32\). By (I2), \(A\) is of type (N), (CM) or (S).

\emph{Type (N).} By (I2), \(A_{\bar\Q}\) contains an elliptic curve \(E\) that is defined over a quadratic field
\(K\). This contradicts \hyperref[thm-quadratic]{Theorem~6}.

\emph{Type (CM).} Here \(A_{\bar\Q}\) is isogenous to \(E^2\) for a CM elliptic curve \(E\), which has potentially
good reduction at \(5\) by (I7)(e). Over a finite extension of \(\Q_5\), \(A\) becomes isogenous to \(E^2\) and
\(E\) acquires good reduction; by (I1)(a), \(A\) acquires good reduction there. So \(A\) has potentially good
reduction at \(5\). The group \(A(\Q)\) contains a subgroup of order \(32\) of \(A(\Q)[2^\infty]\) and a point of
order \(3\), hence a subgroup of order \(96\). This contradicts \hyperref[lemma-five]{Lemma~1}(b).

\emph{Type (S).} By \hyperref[prop-rot]{Proposition~8}, \(\Gamma\cong C_2\) or \(\Gamma\cong V_4\). If
\(\Gamma\cong C_2\), then \(K=L\) is a quadratic field, and by (I2) every elliptic curve in \(A_{\bar\Q}\) is
defined over \(K\); this contradicts \hyperref[thm-quadratic]{Theorem~6}. If \(\Gamma\cong V_4\), we have a
contradiction with \hyperref[prop-v4]{Proposition~9}(c). \qed

\medskip
\emph{Proof of the Corollary.} Let \(A\) be simple over \(\Q\) and geometrically split, and suppose that \(96\)
divides the order of the torsion subgroup of \(A(\Q)\). Then \(A(\Q)\) has a point of order \(3\), and
\(\#A(\Q)[2^\infty]\) is a power of \(2\) divisible by \(32\). This contradicts the Theorem. The Jacobian of a
curve of genus \(2\) over \(\Q\) is an abelian surface, and the groups \([2,48]\), \([2,2,24]\) and
\([2,2,2,12]\) have order \(96\). \qed

\medskip
On the census, \([2,48]\), \([2,2,24]\) and \([2,2,2,12]\) become impossible in the class ``geometrically split,
simple over \(\Q\)''. The groups \([2,48]\) and \([2,2,24]\) have examples split over \(\Q\) (the census curves
\href{https://genus-2-torsion.github.io/curve.html?id=2.48.a}{2.48.a} and
\href{https://genus-2-torsion.github.io/curve.html?id=2.2.24.a}{2.2.24.a}), so their combined geometrically split
cells keep their grades. For \([2,2,2,12]\) the census has only a geometrically simple example
(\href{https://genus-2-torsion.github.io/curve.html?id=2.2.2.12.a}{2.2.2.12.a}); the class ``split over \(\Q\)''
remains open for this group, and with it the combined geometrically split cell.

\section{Remarks}\label{remarks}

\begin{enumerate}
\item\label{rem-where} \emph{Where the hypotheses are used.} Simplicity over \(\Q\) gives \(E^\sigma\neq E\) and
the final contradiction in \hyperref[thm-quadratic]{Theorem~6}; \(\Gamma\neq C_1\) in type (S); and, in type (N),
that the two elliptic curves are defined over a quadratic field and exchanged by \(\tilde\sigma\) ((I2)). It is
needed:
the Jacobians of the census curves 2.48.a and 2.2.24.a are split over \(\Q\), with rational torsion subgroups of
order \(96\). The point of order \(3\) is needed: the census curves
\href{https://genus-2-torsion.github.io/curve.html?id=2.2.8.c}{2.2.8.c},
\href{https://genus-2-torsion.github.io/curve.html?id=2.4.4.c}{2.4.4.c} and
\href{https://genus-2-torsion.github.io/curve.html?id=2.2.2.4.c}{2.2.2.4.c} are geometrically split and simple
over \(\Q\), with rational torsion subgroups \([2,2,8]\), \([2,4,4]\) and \([2,2,2,4]\) of order \(32\). Geometric
splitness is needed: the Jacobian of the census curve 2.2.2.12.a is geometrically simple, with rational torsion
subgroup \([2,2,2,12]\). No polarization is used. \hyperref[lemma-five]{Lemma~1},
\hyperref[lemma-curve]{Lemma~2}, \hyperref[lemma-basechange]{Lemma~3}, \hyperref[lemma-split]{Lemma~4} and
\hyperref[lemma-combine]{Lemma~5} do not involve the surface \(A\) of the Theorem.

\item\label{rem-sharp} \emph{Sharpness.} The census curve
\href{https://genus-2-torsion.github.io/curve.html?id=2.2.2.6.c}{2.2.2.6.c},
\[\begin{aligned}y^2+(x^2+x)y={}&6993x^5-25655829x^4+14380189111x^3\\
&+16095094765893x^2+2757085929534360x+1355999975081200,\end{aligned}\]
has a Jacobian that is geometrically split and simple over \(\Q\), with rational torsion subgroup \([2,2,2,6]\):
by its census record, the Jacobian is isogenous over \(\Q\) to the Weil restriction of an elliptic curve over
\(\Q(\sqrt{-11})\); Part~S of the certificate checks the torsion subgroup and that the Frobenius polynomial at
\(29\) is irreducible, which gives simplicity over \(\Q\). So the bound \(16\) of the Theorem is attained.

\item\label{rem-open} \emph{What the Theorem does not decide.} A rational torsion subgroup of order \(48\) has a
point of order \(3\) and \(2\)-primary part of order \(16\), which the Theorem allows. So the Theorem says nothing
about the groups \([48]\), \([2,24]\), \([4,12]\) and \([2,2,12]\), for which the census has no example in the
class ``geometrically split, simple over \(\Q\)'' on 15 November 2026, nor about \([24]\).

\item\label{rem-two} \emph{Without the point of order \(3\).} \hyperref[prop-rot]{Proposition~8} and
\hyperref[prop-v4]{Proposition~9}(a) do not use a point of order \(3\): if \(A/\Q\) is of type (S), simple over
\(\Q\), and \(\#A(\Q)[2^\infty]\ge32\), then \(\Gamma\cong C_2\), or \(\Gamma\cong V_4\) and \(\Q(i)\subset L\).

\item\label{rem-bn} \emph{Lemma~3 and the literature.} \hyperref[lemma-basechange]{Lemma~3} is the analogue for
\(\Z/2\times\Z/12\) of \cite[Theorem~1.1]{BN}, by which every elliptic curve over a quadratic field with a point of
order \(16\) is the base change of an elliptic curve over \(\Q\). The introduction of \cite{BN} contains a
sentence that begins ``We found examples of elliptic curves that are not base changes of elliptic curves over
\(\Q\) with all possible torsion groups over quadratic fields apart from \(C_{16}\)'' (it continues with the
torsion groups over cubic fields). By Lemma~3 there is no such example with the torsion group
\(\Z/2\times\Z/12\), so this statement is not correct for that group. The sentence reports examples and is not
used in the proofs of \cite{BN}; Theorem~1.1 of \cite{BN} is not affected. Because of this
disagreement, Part~G of the certificate tests Lemma~3 in a second way, which uses neither the involutions
\(\iota_R\) nor the functions \(\lambda_R\): over \(K=\Q(\sqrt{13})\) the curve \(X\) has the point \(P_0\) of
infinite order with \(t=2\sqrt{13}-7\), and for each of the twelve points \(nP_0+T\) (\(1\le n\le3\),
\(T\in X(\Q)\)) the parameter \(t\) is not rational, the curve \(E_t\) has torsion subgroup \(\Z/2\times\Z/12\)
over \(K\), its \(j\)-invariant is rational, and \(E_t\) is isomorphic over \(K\) to the base change of an
elliptic curve over \(\Q\).

Two facts lie behind Lemma~3. First, for each of the four involutions \(\iota_R\), \(R\in X(\Q)\), there is a
function \(\lambda_R\in\Q(X)\) with \(c_4\circ\iota_R=\lambda_R^4c_4\) and \(c_6\circ\iota_R=\lambda_R^6c_6\)
(Part~E). So if \(Q\) is a quadratic point of \(X\), not a cusp, with \(\sigma(Q)=\iota_R(Q)\), then the
invariants of the conjugate of \(E_{t(Q)}\) are \(u^4c_4(Q)\) and \(u^6c_6(Q)\) with \(u=\lambda_R(Q)\in K^\times\),
and \(E_{t(Q)}\) is isomorphic over \(K\) to its conjugate by (I7)(a). (That \(\iota_R\) preserves the
\(j\)-invariant of \(E_t\) shows only that the \(j\)-invariant of \(E_{t(Q)}\) is rational, which is weaker.) This
alone does not make \(E_{t(Q)}\) a base change. Second,
\(\lambda_R\cdot(\lambda_R\circ\iota_R)=1\) (Part~E); this gives \(u\,\sigma(u)=1\) in the proof, which is the
condition for the descent to \(\Q\).

In modular terms (this description is not proved here and it is not used), a point of \(X\) that is not a cusp
is a pair \((E,\varphi)\) up to isomorphism, where \(\varphi\) is an embedding of \(G=\Z/2\times\Z/12\) into
\(E\). The group of automorphisms of \(G\) modulo \(\pm1\) has order \(8\) and acts on \(X\) by
\((E,\varphi)\mapsto(E,\varphi\circ a)\); these are the eight automorphisms \(P\mapsto\pm P+R\) of \(X\) with
\(R\in X(\Q)\). If \(\sigma(Q)=\iota_R(Q)\) and \(\iota_R\) is given by the automorphism \(a\) of \(G\), there is
an isomorphism \(\mu\colon E^\sigma\to E\) over \(K\) with \(\mu\circ\varphi^\sigma=\varphi\circ a\), and
\(\mu\circ\mu^\sigma\) acts on \(\varphi(G)\) as \(a^2\), where \(a^2=\pm1\) because \(\iota_R\) is an involution.
Every automorphism of \(G\) acts as \(\pm1\) on the subgroup of order \(3\), so \(a^2\neq-1\). Hence \(a^2=1\) and
\(\mu\circ\mu^\sigma\) is the identity, which is the condition for the descent of \(E\) to \(\Q\). The proof of
\cite[Theorem~1.1]{BN} has this form, with \(a\) the multiplication by \(9\) on \(\Z/16\), and
\(9^2\equiv1\bmod16\). For the group \(\Z/2\times\Z/10\) there are automorphisms \(a\) with \(a^2=-1\), for
instance the multiplication by \(2\) on the subgroup of order \(5\) together with the identity on the
\(2\)-torsion, and the analogue of Lemma~3 is false. An example (Part~G): over \(\Q(\sqrt{17})\), the curve
\(E(b,c)\) with \(b=(5-3\sqrt{17})/72\) and \(c=(\sqrt{17}-7)/24\) has torsion subgroup \(\Z/2\times\Z/10\), its
\(j\)-invariant \(141420761/9216\) is rational, and it is isomorphic to its conjugate, but it is not the base
change of an elliptic curve over \(\Q\). Indeed, the invariants of a curve isomorphic to \((E_0)_K\) are
\(v^4c_4(E_0)\) and \(v^6c_6(E_0)\) with \(v\in K^\times\), so the norm of \(c_6/c_4\) is the square of a
rational number if \(c_4\neq0\); for this curve \(c_4\neq0\), and the norm of \(c_6/c_4\) is \(17\) times the
square of a nonzero rational number. So for \(\Z/2\times\Z/10\) the statement of \cite{BN} is correct.

\item\label{rem-val} \emph{Validation (not used in the proof).} (i) The script
\texttt{qsimple\_96\_validation.m} takes thirteen curves of the LMFDB \cite{LMFDB} whose Jacobians are of type
(S), with Sato--Tate groups \(E_2\), \(J(E_1)\), \(E_3\), \(E_4\), \(E_6\), \(J(E_2)\), \(J(E_3)\), \(J(E_4)\) and
\(J(E_6)\), which correspond to \(\Gamma=C_2\), \(C_2\), \(C_3\), \(C_4\), \(C_6\), \(V_4\), \(D_6\), \(D_8\) and
\(D_{12}\). For each curve and each prime \(3\le p<500\) of good reduction for the given model, the
\(L\)-polynomial at \(p\) has one of the shapes allowed for that curve: the square of a quadratic polynomial (for a
Frobenius element in \(G_L\)), one of the polynomials of \cite[Lemma~2]{P612} behind (I4), or a polynomial
\((X^2+\varepsilon p)^2-cw^2X^2\) of \hyperref[lemma-inv]{Lemma~7} with \(c\) in a list of at most three squarefree integers
that depends on the curve; every shape allowed for a curve occurs; and the congruences of
\hyperref[prop-rot]{Proposition~8} and \hyperref[prop-v4]{Proposition~9} hold whenever \(32\) or \(3\) divides
the number of points. For the three curves with \(\Gamma=V_4\) the classes \(c\) that occur are \(1,-1\) (once)
and \(2,-1\) (twice); this is consistent with \(c_1c_2c_3\in-\Q^{\times2}\), which then gives the classes
\(1,-1,1\) and \(2,-1,2\) for the three involutions. (ii) The script
\texttt{qsimple\_96\_lmfdb.py} queries the LMFDB tables \texttt{g2c\_curves} (\(66158\) curves) and
\texttt{g2c\_curves\_new} (\(6216959\) curves). Among the curves that are geometrically split and simple over
\(\Q\) and have a rational torsion point of order \(3\) (\(30\) and \(1913\) curves), the \(2\)-primary part of
the rational torsion has order at most \(4\) and at most \(8\) respectively; the four curves with \(2\)-primary
part of order \(8\) have torsion \([2,2,6]\). The curves with Sato--Tate group \(E_n\) or \(J(E_n)\),
\(n\in\{3,4,6\}\), have \(2\)-primary torsion of order at most \(8\), in agreement with
\hyperref[prop-rot]{Proposition~8}. The only curves that are geometrically split and simple over \(\Q\) with
\(32\) dividing the order of the rational torsion are two curves of \texttt{g2c\_curves\_new} with Sato--Tate
group \(N(\mathrm{SU}(2)\times\mathrm{SU}(2))\) (for a geometrically split surface this means type (N)) and torsion of order \(32\). The only curve in
either table with rational torsion of order divisible by \(96\) has torsion \([2,48]\) and a Jacobian that is not
simple over \(\Q\).
\end{enumerate}

\section{Computations}\label{computations}

The certificate
\href{https://genus-2-torsion.github.io/data/knowledge/sources/qsimple_96_checks.m}{\texttt{qsimple\_96\_checks.m}}
(Magma V2.30-1, about \(2\) seconds) consists of assertions and ends with \texttt{ALL\_CHECKS\_PASSED}; a failed
assertion stops it before that line. Its log is
\href{https://genus-2-torsion.github.io/data/knowledge/sources/qsimple_96_checks.log}{\texttt{qsimple\_96\_checks.log}}.

Part~A: \hyperref[lemma-five]{Lemma~1}(a), by enumeration of all polynomials \(h\), and the elliptic curves over
\(\F_7\) with \(8\) and \(12\) points. Part~B: Kubert's family, the exact order of \((0,0)\), the discriminant
and the invariants of \(E_t\); the division polynomials of \(E(b,c)\), their values at \(x=0\), and the inverse
map \(T(b,c)\). Part~C: the \(2\)-division polynomial of \(E_t\) and its discriminant. Part~D: the curve \(X\)
over \(\Q\). Part~E: the involutions \(\iota_R\) and the functions \(\lambda_R\) of
\hyperref[lemma-basechange]{Lemma~3}. Part~F: the three quadratic twists, the reductions at \(13\) and \(37\),
and the \(16\) points of \(X(\Q(\zeta_{12}))\). Part~G: the examples for
\hyperref[lemma-basechange]{Lemma~3}, including those over \(\Q(\sqrt{13})\), and the curve over
\(\Q(\sqrt{17})\) of \hyperref[rem-bn]{Remark~5}. Part~H: the congruences of
\hyperref[prop-rot]{Proposition~8} and \hyperref[prop-v4]{Proposition~9}, and the polynomial identity of
\hyperref[lemma-inv]{Lemma~7}. Part~S: the curve of \hyperref[rem-sharp]{Remark~2}. Every value quoted in the
text from the certificate is asserted, not only printed. Parts~A, B, C, E and H are exact computations with
integers, polynomials and rational functions. Parts~D and~F use Magma's \texttt{CremonaReference} and \texttt{Rank}
(with the flag that the rank is proved) for elliptic curves over \(\Q\), and \texttt{TorsionSubgroup} over \(\Q\)
and, only as a cross-check of \hyperref[lemma-curve]{Lemma~2}(c), over \(\Q(\zeta_{12})\); Part~G uses \texttt{TorsionSubgroup} and \texttt{IsIsomorphic} for elliptic curves over
quadratic fields; Part~S uses \texttt{TorsionSubgroup} and \texttt{LPolynomial} for a curve of genus \(2\). Parts~G
and~S are not used in the proof.

The validation of \hyperref[rem-val]{Remark~6} is not used in the proof. The script
\href{https://genus-2-torsion.github.io/data/knowledge/sources/qsimple_96_validation.m}{\texttt{qsimple\_96\_\allowbreak validation.m}}
(Magma V2.30-1, about \(3\) minutes) ends with \texttt{VAL\_DONE}; its log is
\href{https://genus-2-torsion.github.io/data/knowledge/sources/qsimple_96_validation.log}{\texttt{qsimple\_96\_\allowbreak validation.log}}.
The script \texttt{qsimple\_96\_lmfdb.py} (Python with \texttt{psycopg2}, through the mirror
devmirror.lmfdb.xyz) ends with \texttt{LMFDB\_CHECK\_OK}; its log \texttt{qsimple\_96\_lmfdb.log} was produced on
15 November 2026.

The SHA-256 hashes of the scripts and logs are:
\begin{itemize}
\item \texttt{qsimple\_96\_checks.m}: \texttt{ad292a2a\allowbreak 65541a3a\allowbreak ac660e5f\allowbreak 2e75f0d5\allowbreak 895eca39\allowbreak 9cdaf5ec\allowbreak 45c09b39\allowbreak c507b42f}
\item \texttt{qsimple\_96\_checks.log}: \texttt{cba10806\allowbreak 98adef84\allowbreak 03d4f152\allowbreak b9a0d2ae\allowbreak 3b38cbd6\allowbreak 4cf541d0\allowbreak 593dea04\allowbreak 98ba389a}
\item \texttt{qsimple\_96\_validation.m}: \texttt{79be01f8\allowbreak 6f8c843d\allowbreak 705acc21\allowbreak 22abdaab\allowbreak 0c9b8960\allowbreak 7e4cc05d\allowbreak 61a79cd9\allowbreak 1010261e}
\item \texttt{qsimple\_96\_validation.log}: \texttt{ec47ac83\allowbreak f747dff8\allowbreak 048c9dc2\allowbreak c18c6050\allowbreak 10c4959f\allowbreak d6617464\allowbreak 9637022c\allowbreak 9014bfd5}
\item \texttt{qsimple\_96\_lmfdb.py}: \texttt{3fbba27e\allowbreak 552b4a77\allowbreak bfbd8259\allowbreak 3c2d3f72\allowbreak 2d78cff1\allowbreak 5c12aafa\allowbreak 1ccca17d\allowbreak 6baead15}
\item \texttt{qsimple\_96\_lmfdb.log}: \texttt{ac51cb4a\allowbreak c42d425b\allowbreak 51c30ff2\allowbreak ac5f4012\allowbreak e36a0c1d\allowbreak bc0b1e3a\allowbreak 6c02c701\allowbreak 4aec0b0e}
\end{itemize}

\begin{thebibliography}{LMFDB}

\bibitem[BN]{BN}
P. Bruin and F. Najman, \emph{Fields of definition of elliptic curves with prescribed torsion}, Acta Arith.
\textbf{181} (2017), 85--95: Theorem~1.1, its proof, and the introduction.
\url{https://doi.org/10.4064/aa170323-20-9}, \url{https://arxiv.org/abs/1607.05394}.

\bibitem[Ku]{Ku}
D. S. Kubert, \emph{Universal bounds on the torsion of elliptic curves}, Proc. London Math. Soc. (3)
\textbf{33} (1976), 193--237: Table~3.
\href{https://doi.org/10.1112/plms/s3-33.2.193}{doi:10.1112/plms/s3-33.2.193}.

\bibitem[LMFDB]{LMFDB}
The LMFDB Collaboration, \emph{The L-functions and modular forms database}, \url{https://www.lmfdb.org}, 2026: the
tables \texttt{g2c\_curves} and \texttt{g2c\_curves\_new} (the extended genus-\(2\) data of Booker and
Sutherland), read through the mirror devmirror.lmfdb.xyz.

\bibitem[P40]{P40}
Census proof of 14 November 2026, \emph{Rational \(5\)-torsion and the \(2\)-primary torsion of geometrically split
abelian surfaces}, \url{https://genus-2-torsion.github.io/proofs/qsimple-40-obstruction.html}: Section~1.

\bibitem[P63]{P63}
Census proof of 21 October 2026, \emph{An obstruction to rational \(63\)-torsion on geometrically split abelian
surfaces}, \url{https://genus-2-torsion.github.io/proofs/qsimple-63-obstruction.html}: Sections~2--5 and
Lemmas~2 and~3.

\bibitem[P612]{P612}
Census proof of 15 November 2026, \emph{Rational torsion containing \((\Z/3)^2\) forces \(2\)-primary torsion of
order at most \(4\) on principally polarized abelian surfaces that are geometrically split and simple over
\(\Q\)}, \url{https://genus-2-torsion.github.io/proofs/qsimple-612-obstruction.html}: (I2), (I3) and Lemma~2.

\bibitem[P2210]{P2210}
Census proof of 11 November 2026, \emph{Rational \(5\)-torsion and the \(2\)-rank of geometrically split
Jacobians}, \url{https://genus-2-torsion.github.io/proofs/qsimple-2210-obstruction.html}: (I1).

\bibitem[P2214]{P2214}
Census proof of 10 November 2026, \emph{Rational \(7\)-torsion and the \(2\)-rank of geometrically split abelian
surfaces}, \url{https://genus-2-torsion.github.io/proofs/qsimple-2214-obstruction.html}: Section~1 (with (I7)),
Section~2 and Lemma~3.

\bibitem[Si1]{Si1}
J. H. Silverman, \emph{The Arithmetic of Elliptic Curves}, second edition, Graduate Texts in Mathematics~106,
Springer, Dordrecht, 2009: Section~III.1, Exercise~III.3.7, Proposition~VII.3.1 and Proposition~VII.5.5.

\bibitem[Si2]{Si2}
J. H. Silverman, \emph{Advanced Topics in the Arithmetic of Elliptic Curves}, Graduate Texts in Mathematics~151,
Springer, New York, 1994: Theorem~II.6.1.

\bibitem[ST]{ST}
J.-P. Serre and J. Tate, \emph{Good reduction of abelian varieties}, Ann. of Math. (2) \textbf{88} (1968),
492--517: \S1 (Theorem~1, Corollary~2 and Lemma~2) and \S2.
\href{https://doi.org/10.2307/1970722}{doi:10.2307/1970722}.

\end{thebibliography}

\end{document}
